\chapter{Теоремы о промежуточных значениях непрерывной функции.}
\section{Промежуточные значения непрерывной функции на отрезке}

\begin{thm} [Больцано"--~Коши о промежуточных значениях\rindex{теорема!Больцано"--~Коши о промежуточных значениях}] \label{ch3n2}
Пусть функция~$f$ непрерывна на отрезке~$[a;b]$ и $f(a) = A$,~$f(b) = B$, тогда для любого~$C$, заключенного между~$A$ и~$B$, существует такая точка $c \in [a,b]$, что $f(c)=C$.
\end{thm}
Эту теорему можно переформулировать следующим образом.
\begin{thmn}[Больцано"--~Коши о промежуточных значениях] Непрерывная на отрезке функция, принимая какие-либо два значения, принимает и любое лежащее между ними значение. 
\end{thmn}
\begin{proof}
Пусть для определенности $f(a)=A < B = f(b)$ и тогда $f(a) < C < f(b)$.

Через $c_1$ обозначим середину отрезка $[a;b]$. Если $f(c_1) = C$, то искомая точка найдена и утверждение доказано.

Пусть $f(c_1)\ne C$. Тогда, если $f(c_1) > C$, то положим $a_1 = a$, $b_1 = c_1$, а если $f(c_1) < C$, то $a_1 = c_1$, $b_1 = b$, и поэтому всегда $f(a_1)<C<f(b_1)$.

По тому же самому принципу снова разделим отрезок $[a_1; b_1]$ пополам, и через $c_2$ обозначим его середину. Если $f(c_2) = C$, то утверждение доказано. Если же $f(c_2) \ne C$, то через $[a_2;b_2]$ обозначим ту половину, для которой $f(a_2) < C < f(b_2)$, и т.д. Процесс или обрывается на некотором шаге, и тогда утверждение доказано, или получается последовательность вложенных отрезков $\{[a_n;b_n]\}$ таких, что $\lim_{n \to \infty}\limits(b_n-a_n)=0$ и
\begin{equation}\label{ch3:th1:1}
f(a_n)<C<f(b_n)\quad \forall n\in \bbN.
\end{equation}

Таким образом, согласно  \hyperref[ch1:th:poslstyag]{теореме о последовательности стягивающихся отрезков} существует общая точка $c \in [a;b]$ всех отрезков $\{[a_n;b_n]\}$, причем 
$$
\lim_{n \to \infty}a_n = \lim_{n \to \infty}b_n = c.
$$

Поэтому в силу непрерывности функции~$f$ в точке~$c$:
\begin{equation}\label{ch3:th1:2}
\lim_{n \to \infty}f(a_n) = \lim_{n \to \infty}f(b_n) = f(c),
\end{equation}

Тогда из неравенств~\eqref{ch3:th1:1} в пределе при $n \to \infty$ следует 
\begin{equation}\label{ch3:th1:3}
\lim_{n \to \infty}f(a_n) \le C \le \lim_{n \to \infty}f(b_n),
\end{equation}
откуда из~\eqref{ch3:th1:2} и~\eqref{ch3:th1:3} получаем $f(c)=C$.

Случай, когда $A > B$, рассматривается аналогично.
 
Теорема доказана.
\end{proof}

\begin{cons}
Пусть функция~$f$ непрерывна на отрезке~$[a;b]$ и $f(a)$, $f(b)$ имеют разные знаки, тогда существует такая точка $c \in [a,b]$, что $f(c)=0$.
\end{cons}

\begin{thm}
Пусть функция $f$ непрерывна на отрезке~$[a;b]$ и $M=\sup f([a;b]) $, $m=\inf f([a;b])$. Тогда функция $f$ принимает все значения из отрезка $[m;M]$ и только эти значения.
\end{thm}
\begin{proof}
Согласно теореме~\ref{th:ch2:Veyershtrass} (Вейерштрасса) из предыдущей главы существуют такие точки $\alpha\in [a;b]$ и $\beta\in [a;b]$, что $f(\alpha)=m$, $f(\beta)=M$

Тогда рассматриваемая теорема непосредственно вытекает из теоремы~\ref{ch3n2} (Больцано"--~Коши), примененной к отрезку $[\alpha;\beta]$, если $\alpha \le \beta$, или соответственно к отрезку $[\beta;\alpha]$, если $\beta \le \alpha$. 

Теорема доказана.
\end{proof} 

Таким образом, множество всех значений функции, заданной и непрерывной на некотором отрезке, представляет собой также отрезок. Однако, как показывают примеры, образом интервала может быть любой промежуток.

\section{Про обратное утверждение}
Заметим, что обратное утверждение является неверным, например, $f(x)= \sin\frac{1}{x}$ для $x\ne0$ и $f(0) = 0$ разрывна в точке $x = 0$, но у нее образ любого отрезка есть отрезок. Однако, для монотонных функций обратное утверждение является верным, докажем это в Теореме~\ref{ch3n3}.

